\documentclass[]{report}


% Title Page
\title{}
\author{}
\usepackage{amsmath}
\usepackage{float}
\usepackage{hyperref} 
\usepackage{graphicx}

\usepackage[space]{grffile}

\begin{document}
\maketitle

\begin{abstract}

\end{abstract}

\chapter{Introduction}
\section{Introduction}
My goal in this thesis is to study and better understand the dynamics of ellipsoidal particles in shear flows. The thesis is a continuation of two previous MSc theses\cite{AntonThesis, JonasThesis}. The methodology is to experimentally measure the orientational dynamics of micrometer length glass particles in a shear flow and comparing the results to those of theoretical models. In the first part of the thesis I describe the improvements that were made to the experimental setup, most importantly an automated tracking. In the second part of the thesis the measurements and their analysis is discussed. But before discussing either of these subjects more in depth some background and theory is needed.

\subsection{Background}
Understanding the orientational dynamics of particles in flow might appear somewhat esoteric to someone unfamiliar with the field, but there is a number of topics where it is very useful. In medical applications understanding the dynamics of ellipsoidal particles such as bacteria can be relevant to a detailed understanding of their interactions with cells and other bodies. This is discussed by Tolga \emph{et al}~\cite{Tolga}. 

One of the most influential papers in the study of particle dynamics in flow was by Einstein in 1905~\cite{Einstein}. He showed how much suspended spherical particles would increase the viscosity of a fluid. Jeffery in his 1922 paper~\cite{Jeffery} extended these results to ellipsoidal  particles and derived equations for the orientational dynamics of axisymmetric particles, in other words how the particles would rotate as a function of time. For systems where inertial effects could be disregarded the motion was found to be periodic and depending only on the initial condition of the particle. 

Investigation of triaxial particles was started by Gierszewski \& Chaffey~\cite{Chaffey} and was continued by Hinch \& Leal~\cite{Leal} and more recently by Yarin \emph{et al}~\cite{Yarin}. 
The dynamics Jeffery had found for axisymmetric particles were periodic, but it was shown by Hinch \& Leal that for triaxial particles some orbits would be doubly periodic, in other words following two separate independent periods. This behaviour will in this thesis be referred to as \emph{quasi-periodic}.

Yarin \emph{et al} used numerical simulations to generate a surface of section~\cite{SurfaceOfSection} for ellipsoidal particles with different shapes. They showed that not only were there double periodic or quasi periodic orbits but when the particles were sufficiently different from axisymmetric there would be chaotic orbits. 
Several other surfaces of section were produced by Johansson ~\cite{AntonThesis} using the same method as Yarin. It was shown that even small asymmetries of the order of 1\% lead to quasi-periodic motion for some initial conditions.

Attempts to experimentally verify these theoretical results were initially performed by Goldsmith and Mason in 1962~\cite{Mason} who used flow in a glass pipe to observe the rotation rate for several different particle shapes. They confirmed that the rotation rate matched well with that predicted from Jeffery orbits but they did not study the actual orbits. Since then most experimental research, such that as by Harlen and Koch~\cite{fibersspension} has focused on how diluted suspensions of particles affect the properties of a liquid. Only tangential efforts such as by Tolga~\cite{Tolga} were concerned with the Jeffery orbits. A good summary of both theoretical and experimental results was written by Petrie~\cite{Petrie} in 1999.

The first dedicated experiments to measure the actual Jeffery orbits in angular components and verify the orientational dynamics were performed by Einarsson \emph{et al}~\cite{JonasExperiment}. Although there were some promising results, the vast majority of particles were asymmetric to the degree that their orbits were chaotic or highly quasi-periodic. Moreover the width and length of particles varied greatly and could not be measured accurately.
This meant that although the orbits could be qualitatively shown to be similar to some 
Jeffery orbits, no particular particle could be shown to exhibit both quasi periodic and periodic motion. No particle could also be well matched to a particular orbit, but different particles could be shown to be qualitatively simililar to different types of orbits. There was also few particles that very closely retraced its trajectory along the entire length of the channel when the flow was reversed, which indicates that their deviation from periodic motion was not caused by noise, as that would not be reversible.


% I really want a cite for thus but how could I possibly do that.
The goal of this thesis is to experimentally verify the results of Yarin and Hinch, Leal\cite{Yarin, Leal} and show that the same particle will show different types of motion for different initial conditions. Furthermore that different particles will show different motion for the same initial conditions based on the asymmetry. This is done by observing the orientation of a micrometer length particle in a creeping shear flow. The flow is shown to be creeping by demonstrating that the particle dynamics revert as the flow is reverted. The results are then will then be compared to theoretical predictions for different initial conditions and asymmetries.

%doing so would be very important to actually motivating using these results as well as possibly finding the limitations of this theory in real world applications. Understanding the dynamics of ellipsoidal particles in shear flow could be useful for example in our understanding of microscopic bacteria in blood stream 

\section{Background}
	

\section{Experimental Setup}
In order to study the dynamics of particles, a channel made from PDMS and bonded onto a glass slide. A fluid containing the traxial glass particles that can be seen in fig re REF HERE is pumped at a constant rate through the channel using a INSERT THE PUMP NAME HERE. To observe the particles, the channel is placed upon a confocal microscope connected to a INSERT CAMERA MODEL HERE. The channel is then moved using a INSERT STEP ENINE NAME HERE controlled from a computer. A mroe detailed description of the setup can be seen in figure REF to SETUP IMAGE THAT I MAKE

\subsection{The particles}
	Previously the particles used in the experiment were polymeric made by swirling an apoxy solution in a vortex. This produces rod like particles of different lengths, but very often included noticeable defects such as being bent, small lumps etc. 
	To improve results glass particles from Nippon Glass company in Japan were brought in. As they are mainly manufactured to serve as spacing rods in LCD screens they have a very well determined width but then broken at varying lengths. These breaking points are the only real visible point of imperfection of the rods.
	The specification of length and width can be seen in figure \ref{fig:rod_lengths}
	
	While the glass rods are far more regular in terms of size and shape, they do have on major problem. The density of the glass is 2.57 $g/cm^3$, significantly higher than water which has a density of around ENSITY OF WATER at 20 degrees C and glycerol which has a density of DENSITY OF GLYCEROL. And so the particles will sink very quickly in any mix of water and glycerol.
	
	The solution is to use Sodium Metatungstate, water soluble mineral which gives the resulting mix a maximum density of around DENSITY OF SOMT which is more than sufficient to keep the particles suspended. It is however not without its own issues which will be discussed in the discussion portion of the thesis. 
	
	
	First the viscosity of the mix is significantly lower than the glycerol and water mix used previously (a 1:4 mix of glycerol and water was commonly used) which means the system is more sensitive to noise. 
	Second the density of the solution changes when the water mixed in evaporates. As the ratio of water to mineral is relatively high, clearly more than 1:1 by weight, a small percentage of water evaporating will result in significant change in density
	
	
	
	
\subsection{Sodium Metatungstate}
I guess we use sodium meta tungstate. But does this really need a whole section? Seems like overdoing it... but what %else have I done v_v
THE EQUATION FOR MODIFYING VISCOCITY

\begin{equation}
\rho_{a} = 	\frac {m_{a}}{V_{a}} =
				\frac{V_{b}\rho_{b} + V_{mix}\rho_{mix}}{V_b + V_{mix}} 
\end{equation}
So if we want to find $V_{mix}$ we get
\begin{equation}
V_{mix} = \frac{ V_{b}(\rho_{b} - \rho_{a})}{\rho_{a} - \rho_{mix}} 
\end{equation}



\chapter{Theory}
In order to better understand the results and discussions in this thesis some rather basic results in fluid dynamics, image analysis and other fields need to be discussed. It is however not the focus of this thesis and as such I will try to be as brief as possible, and an experienced reader can safely skip to the next chapter.


\section{Fluid Dynamics}
In order to understand the limitations and behaviours of the experiment we need to know about a few key concepts in fluid dynamics.
\subsection{Reynold's Number}
The Reynolds number (Re) is a dimensionless number describing the ratio of intertial forces to viscous forces in a flow. This is not a very strict definition, but suffices to explain that the Reynolds number can be used to characterize the so called flow regime of a system. 

There are two major types of flows, Laminar flow and turbulent flow, sometimes also talking about transitional flow being the region between the two but it is not of great importance to this thesis. 

Laminar flow. nonturbulent streamline flow in parallel layers. There are no cross currents perpendicular to the direction of the flow, which in this experiment is relatively straight forward. An example of a laminar flow profile can be seen in figure \ref{fig:laminar_flow} as well as turbulent flow. Essentially laminar flow is the regime where the viscous forces are larger than the intertial forces.

\begin{figure}\label{fig:laminar_flow}
\includegraphics{Images/laminarFlow.png}
\end{figure}

Turbulent flow is characterized by violent, unsteady and unpredictable flow. This is the regime where the intertial forces are larger than the viscous forces.

The Reynolds number is the primary tool used to describe the regime of a flow, and while there is no clear number at which a flow stops being laminar but the value usually used is that the flow is laminar if $Re << 1$
\subsection{Stokes Drag and Stokes law}
Stokes Law:
\begin{equation}
v_s = \frac{2}{9} \frac{\delta \rho}{\mu} g R^2
\end{equation}

$\delta \rho$ = 0.01 ... 0.001

$\mu = 8*10^-3$

$g = 9.82$

$R = 5 - 35 \cdot 10^-6$


\section{Jeffrey Orbits}

\section{Kalman filter}
% Should I even have this section?
\chapter{Automated Tracking}
TODO:
READ UP ON THE EFFICIENCY OF THE CANNY EDGE DETECTION
TAKE AN IMAGE OF THE STATIC NOISE REDUCTION (DOES IT ACTUALLY WORK?!)
GET AN IMAGE OF BEFORE AND AFTER SMOOTHING, PREFERABLY INCLUDE A CANNY EDGE DETECTION OF BOTH 
ADD ANTONS THESSIS TO BIBLIOGRAPHY
WORK MORE ON KALMAN FILTER SECTION
WRITE ABOUT THE TIME IT TAKES TO DO VARIOUS STUFFS

\section{Noise Reduction}
The first step in tracking a particle is to correctly identify it in the image given. To do this we want to eliminate as much noise from the image as we can. The first type of noise we will eliminate is static noise, this is noise present in every image and results from dirt or scratches on the lens of the camera or microscope. 
The easiest solution to such noise would of course be to clean the instruments, but despite numerous attempts of carefully cleaning every surface a noticeable of static noise would always remain. 

This means one must use algorithmic noise reduction methods to remove this noise. The method employed is a simple averaging that takes N pictures at different positions in the channel to generate an average image where any particular features of the channel would disappear and only the static noise remain. This can be seen in figure STATIC NOISE REDUCTION IMAGE,

The second form of noise eliminated are objects of no interst, such as the lines in the PDMS channel, from the scratches made polishing the cast the channel is molded from. Or, possible dirt on the cover glass. To reduce this noise, the image is blurred with a gaussian blur filter of size 7 with $\sigma=1$. THE RESULTS OF SMOOTHING CAN BE SEEN IN ANOTHER IMAGE

\section{Canny Edge Detection}
The Canny Edge invented by John Canny in his 2011 paper the canny of cane \cite{Canny}detection is generally considered the most advanced and best performing edge detection of the simple filter functions. 
Without going in to the finer details, given an image matrix $\mathbf{I}$ where each value corresponds to the light intensity $i_{x,y}$of that pixel at index $x,y$ the Canny edge detector will try to find the cohesive pixels $E = \{e_1, e_2... e_n\}$ where there is a noticeable change in intensity. In other words, what we call an "edge" an image. 

This is accomplished in two steps. First a Sobel Filter edge image is computed by looking at the change in intensity at every pixel from 3 possible directions and averaging these. 

S = MATHEMATICS OF SOBEL

We then consider a Then, a pixel 

$p_{x,y} \in E \text{if} S(x,y) > T_{high}$

where $T_{high}$ is a predetermined threshold value. 

Secondly we recursively check all pixels neighbouring an edge

MATH?

 if they are higher than some threshold value $T_low$. if they are, they are also considered part of the edge. This is repeated until no more pixels are added. An image illustrating this process can be seen in figure FIG. The benefit of the Canny edge detection over the simpler Sobel is that it is much easier to detect cohesive objects that vary in intensity without getting a lot of other noise in the image. The only real drawback of the Canny Edge detection is the computation time, and thus the implementation from the Open Computer Vision (OCV) was used, as this is a heavily optimized routine written in C++ with real time uses in mind. Thanks to this the computing for the Canny Edge of a 260x260 Image takes about X ms and thus is of little importance in the overall time per frame.

After noise reduction the Canny Edge Detection is used to find the most significant edges in the image. 

\section{Contour detection and selection}

Once an edge image has been generated, the OCV package has another useful function, \texttt{Contours} which returns a list of every contiguous group of edge pixels. If we have chosen the threshold values to the edge detection correctly, this should include the particle or a good approximation of it. 

In order to find the correct contour, a few techniques are used to find the correct contour.

First particles whose total size is less than some minimum value, $ n_{min}$ or larger than some maximum value $n_{max}$ are ignored. Then the position $pP_i$ of each contour $C_i={p_1,p_2...p_n}$ is calculated as the average pixel position

\[
P_i = \sum_{j}^n p_j/n
\]
This position is compared to the expected position of the Kalman filter, which the very first frame is the middle position. 

Finally a 'thinness value' is calculated according to eq \ref{eq:thinness}

\begin{equation}\label{eq:thinness}
w_{thin}\left(\frac{ n}{d_{max}^2}\right)^2
\end{equation}. 
where $w_{thin}$ is a weighting constant, $n$ is the number of pixels in the contour and $d_{max}$ is the longest distance between two pixels in the 
where I am not really sure I should do this now that I have so few particles, but I do it none the less!



\section{Kalman Filter}
	When the particle is at constant motion in the channel the equations of motion give us
	\begin{equation}
	\left[
	\begin{array}{c}
	  x_n	 \\
	 	  v_{x,n} \\ 
	  y_n	 \\
	  v_{y,n} 
	\end{array} 
	\right]
	 =
	 \begin{bmatrix}
	  x_{n-1} & + v_{x,n-1} \\
	  v_{x,n-1}       & ...		 \\
	  y_{n-1} &+v_{y,n-1}  \\
	  v_{y,n-1}       &		\\
	 \end{bmatrix}
	+
	\left[
	\begin{array}{c}
	0 \\
	c_{xn} \\
	0 \\
	c_{yn}
	\end{array}
	\right]
	%	x_n \right] = \left[ x_{n-1}+ v_{n+1} \right]\\
	%	v_{xn}
		y_n
		v_{yn}
	\end{equation}
	
	We can rewrite this in matrix form as
	
	\begin{equation}
	\left[
	 \begin{array}{c}
	 x_{n} \\
	 v_{x,n}\\
	 y_{n}\\
	 v_{y,n}
	 \end{array}
	 \right]
	 =
	 \begin{bmatrix}
	  1 & 	\Delta t & 	0 & 	0 \\
	  0 & 	1 & 	0 & 		0 \\
	  0 & 	0 & 	1 & \Delta t  \\
	  0 & 	0 & 	0 & 		1 \\
	 \end{bmatrix}
	 \cdot
	 \left[
	 \begin{array}{c}
	 x_{n-1} \\
	 v_{x,n-1}\\
	 y_{n-1}\\
	 v_{y,n-1}
	 \end{array}
	\right]
	 +
	 \begin{bmatrix}
	   0 & 	0 & 	0 & 	0 \\
	   0 & 	1 &		0 & 	0 \\
	   0 & 	0 & 	0 &    	0 \\
	   0 & 	0 & 	0 & 	1 \\
	  \end{bmatrix}
	  \cdot
	  \begin{array}{c}
	  0 \\
	  c_{x,n-1} \\
	  0 \\
	  c_{y,n-1}
	  \end{array}
	\end{equation}
	
	Now if we want to transform this to the measured data we have to simply multiply by a constant pixel-to-meter constant, which has been meassured to be around $1/940=1.06\cdot 10^{-3}$
	
	\begin{equation}
	Q = 
	 \begin{bmatrix}
	  0 & 	0 & 	0 & 	0 \\
	   0 & 	1 &		0 & 	0 \\
	   0 & 	0 & 	0 &    	0 \\
	   0 & 	0 & 	0 & 	1 \\
	  \end{bmatrix}
	\end{equation}
	
	Which we can easily re-write as a Kalman filter 
	\begin{equation}
	\hat{\mathbf{x}}_{n} = \mathbf{A}_n\mathbf{x_{n-1}} + \mathbf{B}_n \mathbf{u}_n + \mathbf{w}_n
	\end{equation}
	\begin{equation}
	\hat{\mathbf{P}}_n = \mathbf{F}_n \mathbf{P}_{n-1} \mathbf{F}^{T}_{n} + \mathbf{Q}_n
	\end{equation}
	with 
	\begin{eqnarray}
	\mathbf{x}_n &= 
	\left[
	x_n 
	v_{x,n} 
	y_n 
	v_{y,n}
	\right]^T \\
	\mathbf{A} &= 
	\begin{bmatrix}
	  1 & 	\Delta t & 	0 & 	0 \\
	  0 & 	1 & 	0 & 		0 \\
	  0 & 	0 & 	1 & \Delta t  \\
	  0 & 	0 & 	0 & 		1 \\
	\end{bmatrix} \\
	\mathbf{B}_n &= 
	 \begin{bmatrix}
	   0 & 	0 & 	0 & 	0 \\
	   0 & 	1 &		0 & 	0 \\
	   0 & 	0 & 	0 &    	0 \\
	   0 & 	0 & 	0 & 	1 \\
	  \end{bmatrix}
	\\
	\mathbf{u}_n &=
	\left[
	  \begin{array}{c}
	  0 \\
	  c_{x,n-1} \\
	  0 \\
	  c_{y,n-1}
	  \end{array}
	  \right] 
	\end{eqnarray}
	
	and we add the update part as
	\begin{eqnarray}
	&\mathbf{K}_n = \hat{\mathbf{P}}_n \mathbf{H}^T_n(\mathbf{H}_n \hat{\mathbf{P}}_n \mathbf{H}^T_n	) \\
	&\mathbf{x_n} = \hat{\mathbf{x}}_n + \mathbf{K}_n(\mathbf{y} - \mathbf{H}_n\hat{\mathbf{x}}_n) \\
	&\mathbf{P}   = (\mathbf{I} - \mathbf{K}_n\mathbf{H}_n)\hat{\mathbf{P}}_n
	\end{eqnarray}
	
	
	
	
	

\section{Stabilizing the tracking}
Once a state estimation has been made, we want to adjust the speed of the step engine to, as best as possible, match that of the particle. This is done by looking both the position and velocity of the particle and going through the conditional statements shown in figure CONDITIONAL CORRECTION VECTOR. 

The goal is to limit the amount of corrections made, as changing the velocity is rather time intensive as discussed in section \ref{sec:time considerations}, as well as keeping the particle stable. There is also no point in trying to completely eliminate movement
% % About how to choose the correction vector and communicate with the step engine

\section{Time Considerations}\label{sec:time considerations}
A higher FPS will allow the particle detection to be better, improve the position saving and allow the particle tracking to be more stable as well. So maximizing the FPS, ie reducing the computational time of each task, is a clear goal for a good automated tracking. A list of the different tasks and their average execution times can be seen in table \ref{tab:benchmarks}

\begin{table}[H]
 \begin{tabular}{l | c | c } 
 Task  			&  Average time & Std deviation \\
 Capture screen & 1000 			& 200 \\
 Find edges 	& 200			& 20 \\
 \end{tabular}\label{tab:benchmarks}
\end{table}

We can clearly see that FPS is limited primarily by three routines: The screen capture routine, the change velocity routine and finally the save position routine. The first and last are unavoidable and must be done every frame by definition if we are interested in knowing the particles position as well as possible. This means we simply want to use the velocity correction as little as possible. Since the time constraint is in the communication with the step engine, there is not any optimization to be done here, at least not within the scope is this thesis. 


\chapter{Experimental Setup}


\chapter{Results}
% LOLOLOLOLOLOL
\chapter{Discussion}

% % STUFF I SHOULD CITE % %
%http://ncds.technion.ac.il/userfiles/publications/jfm97ygr_ellipsoid.pdf
\appendix
\chapter{Raw data}

%\section{$\mathbf{\lambda}$ of about 6}
%These are the particles with a $\lambda = 6 \pm 0.5$
%\subsection{Particle 1}

\begin{figure}[H]
\centering
\includegraphics[width=0.8\textwidth]{Images/Particle 1/Particle1.pdf}
\caption{Particle 1: $ \lambda: 6.1591$Depth: 160 out of $200 \mu $ m}
\end{figure}

\begin{figure}[H]

\centering

\includegraphics[width=0.8\textwidth]{Images/Particle 1/Stretch1.pdf}

\end{figure}

\begin{figure}[H]

\centering

\includegraphics[width=0.8\textwidth]{Images/Particle 1/Stretch2.pdf}

\end{figure}

\subsection{Particle 4}

\begin{figure}[H]
\centering
\includegraphics[width=0.8\textwidth]{Images/Particle 4/Particle4.pdf}
\caption{Particle 4: At around 10 minutes, pretty big "jump" in channel flow. $ \lambda: 12.7745$Depth: 140 out of $200 \mu $ m}
\end{figure}

\begin{figure}[H]
\centering
\includegraphics[width=0.8\textwidth]{Images/Particle 4/Stretch1.pdf}
\end{figure}


\begin{figure}[H]
\centering
\includegraphics[width=0.8\textwidth]{Images/Particle 4/Stretch2.pdf}
\end{figure}


\subsection{Particle 11}

\begin{figure}[H]
\centering
\includegraphics[width=0.8\textwidth]{Images/Particle 11/Particle11.pdf}
\caption{Particle 11: At around 7:03, possible interaction with dirt in channel. $ \lambda: 5.9638$Depth: 180 out of $200 \mu $ m}
\end{figure}

\begin{figure}[H]
\centering
\includegraphics[width=0.8\textwidth]{Images/Particle 11/Stretch1.pdf}
\end{figure}

\begin{figure}[H]
\centering
\includegraphics[width=0.8\textwidth]{Images/Particle 11/Stretch2.pdf}
\end{figure}


\subsection{Particle 13}
\begin{figure}[H]
\centering
\includegraphics[width=0.8\textwidth]{Images/Particle 13/Particle13.pdf}
\caption{Particle 13: At around 4:15 there is a big jump in the flow at reversal. $ \lambda: 5.5457$Depth: 140 out of $200 \mu $ m}
\end{figure}

\begin{figure}[H]
\centering
\includegraphics[width=0.8\textwidth]{Images/Particle 13/Stretch1.pdf}
\end{figure}


\subsection{Particle 16}
\begin{figure}[H]
\centering
\includegraphics[width=0.8\textwidth]{Images/Particle 16/Particle16.pdf}
\caption{Particle 16: $ \lambda: 6.2917$Depth: 120 out of $200 \mu $ m}
\end{figure}

\begin{figure}[H]

\centering

\includegraphics[width=0.8\textwidth]{Images/Particle 16/Stretch1.pdf}

\end{figure}

\begin{figure}[H]

\centering

\includegraphics[width=0.8\textwidth]{Images/Particle 16/Stretch2.pdf}

\end{figure}



\subsection{Particle 22}

\begin{figure}[H]


\centering
\includegraphics[width=0.8\textwidth]{Images/Particle 22/Particle22.pdf}
\caption{Particle 22: $ \lambda: 11.6767$Depth: 160 out of $200 \mu $ m}

\end{figure}

\begin{figure}[H]

\centering

\includegraphics[width=0.8\textwidth]{Images/Particle 22/Stretch1.pdf}

\end{figure}

\begin{figure}[H]

\centering

\includegraphics[width=0.8\textwidth]{Images/Particle 22/Stretch2.pdf}

\end{figure}

\subsection{Particle 23}

\begin{figure}[H]
\centering
\includegraphics[width=0.8\textwidth]{Images/Particle 23/Particle23.pdf}
\caption{Particle 23: $ \lambda: 5.6106$Depth: 80 out of $200 \mu $ m}
\end{figure}

\begin{figure}[H]
\centering
\includegraphics[width=0.8\textwidth]{Images/Particle 23/Stretch1.pdf}
\end{figure}




%
%\section{$\mathbf{\lambda}$ of about 6}
%These are the particles with a $\lambda = 7 \pm 0.5$
%\subsection{Particle 2}

\begin{figure}[ H]

\caption{Particle 2: $ \lambda: 13.9227$Depth: 30 out of $200 \mu $ m}

\centering

\includegraphics[width=0.8\textwidth]{Images/Particle 2/Particle2.pdf}

\end{figure}

\begin{figure}[ H]

\centering

\includegraphics[width=0.8\textwidth]{Images/Particle 2/Stretch1.pdf}

\end{figure}

\begin{figure}[ H]

\centering

\includegraphics[width=0.8\textwidth]{Images/Particle 2/Stretch2.pdf}

\end{figure}

\begin{figure}[ H]

\centering

\includegraphics[width=0.8\textwidth]{Images/Particle 2/Stretch3.pdf}

\end{figure}

\begin{figure}[ H]

\centering

\includegraphics[width=0.8\textwidth]{Images/Particle 2/Stretch4.pdf}

\end{figure}

\begin{figure}[ H]

\centering

\includegraphics[width=0.8\textwidth]{Images/Particle 2/Stretch5.pdf}

\end{figure}

\begin{figure}[ H]

\centering

\includegraphics[width=0.8\textwidth]{Images/Particle 2/Stretch6.pdf}

\end{figure}


\subsection{Particle 7}

\begin{figure}[ H]

\caption{Particle 7: $ \lambda: 27.4694$Depth: 20 out of $200 \mu $ m}

\centering

\includegraphics[width=0.8\textwidth]{Images/Particle 7/Particle7.pdf}

\end{figure}


\subsection{Particle 18}

\begin{figure}[ H]

\caption{Particle 18: Unclear. There are several possible interactions around 6 minutes but neither is all too clear. Also the second turn is overall quite jumpy$ \lambda: 14.3568$Depth: 20 out of $200 \mu $ m}

\centering

\includegraphics[width=0.8\textwidth]{Images/Particle 18/Particle18.pdf}

\end{figure}

\begin{figure}[ H]

\centering

\includegraphics[width=0.8\textwidth]{Images/Particle 18/Stretch1.pdf}

\end{figure}

\begin{figure}[ H]

\centering

\includegraphics[width=0.8\textwidth]{Images/Particle 18/Stretch2.pdf}

\end{figure}


\subsection{Particle 19}

\begin{figure}[ H]

\caption{Particle 19: At around 0:50 there is an uneven vertical flow "surging" the particle down. $ \lambda: 13.795$Depth: 120 out of $200 \mu $ m}

\centering

\includegraphics[width=0.8\textwidth]{Images/Particle 19/Particle19.pdf}

\end{figure}

\begin{figure}[ H]

\centering

\includegraphics[width=0.8\textwidth]{Images/Particle 19/Stretch1.pdf}

\end{figure}

\begin{figure}[ H]

\centering

\includegraphics[width=0.8\textwidth]{Images/Particle 19/Stretch2.pdf}

\end{figure}



 


\subsection{Particle 20}

\begin{figure}[ H]

\caption{Particle 20: At 5:22 collides with large "debris" on channel floor. $ \lambda: 13.5864$Depth: 150 out of $200 \mu $ m}

\centering

\includegraphics[width=0.8\textwidth]{Images/Particle 20/Particle20.pdf}

\end{figure}

\begin{figure}[ H]

\centering

\includegraphics[width=0.8\textwidth]{Images/Particle 20/Stretch1.pdf}

\end{figure}



\subsection{Particle 24}

\begin{figure}[ H]

\caption{Particle 24: $ \lambda: 14.3715$Depth: 60 out of $200 \mu $ m}

\centering

\includegraphics[width=0.8\textwidth]{Images/Particle 24/Particle24.pdf}

\end{figure}

\begin{figure}[ H]

\centering

\includegraphics[width=0.8\textwidth]{Images/Particle 24/Stretch1.pdf}

\end{figure}

\begin{figure}[ H]

\centering

\includegraphics[width=0.8\textwidth]{Images/Particle 24/Stretch2.pdf}

\end{figure}

%
%\section{Other particles}
%\subsection{Particle 3}

\begin{figure}[H]

\includegraphics[width=0.8\textwidth]{Images/Particle 3/Particle3.pdf}

\caption{Particle 3:  $ \lambda: 19.5489$ Depth: 140 out of $200 \mu $ m}

\centering

\end{figure}

\begin{figure}[H]

\centering

\includegraphics[width=0.8\textwidth]{Images/Particle 3/Stretch1.pdf}

\end{figure}


\subsection{Particle 5}

\begin{figure}[H]

\includegraphics[width=0.8\textwidth]{Images/Particle 5/Particle5.pdf}

\caption{Particle 5:  $ \lambda: 15.2445$ Depth: 160 out of $200 \mu $ m}

\centering

\end{figure}

\begin{figure}[H]

\centering

\includegraphics[width=0.8\textwidth]{Images/Particle 5/Stretch1.pdf}

\end{figure}

\begin{figure}[H]

\centering

\includegraphics[width=0.8\textwidth]{Images/Particle 5/Stretch2.pdf}

\end{figure}


\subsection{Particle 6}

\begin{figure}[H]
\includegraphics[width=0.8\textwidth]{Images/Particle 6/Particle6.pdf}
\caption{Particle 6:  $ \lambda: 17.3351$ Depth: 150 out of $200 \mu $ m}
\centering

\end{figure}

\begin{figure}[H]

\centering

\includegraphics[width=0.8\textwidth]{Images/Particle 6/Stretch1.pdf}

\end{figure}

\begin{figure}[H]

\centering

\includegraphics[width=0.8\textwidth]{Images/Particle 6/Stretch2.pdf}

\end{figure}

\subsection{Particle 8}

\begin{figure}[H]

\includegraphics[width=0.8\textwidth]{Images/Particle 8/Particle8.pdf}

\caption{Particle 8:  $ \lambda: 36.1074$ Depth: 120 out of $200 \mu $ m}

\centering

\end{figure}

\subsection{Particle 9}

\begin{figure}[H]

\includegraphics[width=0.8\textwidth]{Images/Particle 9/Particle9.pdf}

\caption{Particle 9:  $ \lambda: 7.5902$ Depth: 160 out of $200 \mu $ m}

\centering

\end{figure}

\subsection{Particle 10}

\begin{figure}[H]

\includegraphics[width=0.8\textwidth]{Images/Particle 10/Particle10.pdf}

\caption{Particle 10:  $ \lambda: 8.6589$ Depth: 60 out of $200 \mu $ m}

\centering

\end{figure}

\begin{figure}[H]

\centering

\includegraphics[width=0.8\textwidth]{Images/Particle 10/Stretch1.pdf}

\end{figure}

\subsection{Particle 12}

\begin{figure}[H]

\includegraphics[width=0.8\textwidth]{Images/Particle 12/Particle12.pdf}

\caption{Particle 12:  $ \lambda: 8.1961$ Depth: 180 out of $200 \mu $ m}

\centering

\end{figure}

\begin{figure}[H]

\centering

\includegraphics[width=0.8\textwidth]{Images/Particle 12/Stretch1.pdf}

\end{figure}


\subsection{Particle 15}

\begin{figure}[H]

\includegraphics[width=0.8\textwidth]{Images/Particle 15/Particle15.pdf}

\caption{Particle 15:  $ \lambda: 7.7306$ Depth: 120 out of $200 \mu $ m}

\centering

\end{figure}

\begin{figure}[H]

\centering

\includegraphics[width=0.8\textwidth]{Images/Particle 15/Stretch1.pdf}

\end{figure}

\begin{figure}[H]

\centering

\includegraphics[width=0.8\textwidth]{Images/Particle 15/Stretch2.pdf}

\end{figure}


\subsection{Particle 17}

\begin{figure}[H]

\includegraphics[width=0.8\textwidth]{Images/Particle 17/Particle17.pdf}

\caption{Particle 17:  $ \lambda: 7.6562$ Depth: 180 out of $200 \mu $ m}

\centering

\end{figure}

\begin{figure}[H]

\centering

\includegraphics[width=0.8\textwidth]{Images/Particle 17/Stretch1.pdf}

\end{figure}

\begin{figure}[H]

\centering

\includegraphics[width=0.8\textwidth]{Images/Particle 17/Stretch2.pdf}

\end{figure}

\subsection{Particle 21}

\begin{figure}[H]

\includegraphics[width=0.8\textwidth]{Images/Particle 21/Particle21.pdf}

\caption{Particle 21:  $ \lambda: 9.6776$ Depth: 60 out of $200 \mu $ m}

\centering

\end{figure}

\begin{figure}[H]

\centering

\includegraphics[width=0.8\textwidth]{Images/Particle 21/Stretch1.pdf}

\end{figure}

\begin{figure}[H]

\centering

\includegraphics[width=0.8\textwidth]{Images/Particle 21/Stretch2.pdf}

\end{figure}

\bibliography{bibliography}
\bibliographystyle{plain}
\end{document}          
